%% ---------------------------------------------------------------------------
%% conclusiones.tex
%% Conclusiones y recomendaciones
%% ---------------------------------------------------------------------------

\chapter{Conclusiones y recomendaciones}
\label{ch:conclusiones}

\section{Conclusiones}

El presente proyecto logró diseñar, implementar e integrar un módulo PLPUF
programable dentro de un SoC basado en FPGA, cumpliendo satisfactoriamente
el objetivo general propuesto. El módulo opera como periférico AXI4-Lite
controlado por el procesador suave MicroBlaze~V sobre la plataforma Basys\,3
con FPGA Artix-7, proporcionando una solución funcional, abierta y
reproducible para la generación de identificadores de hardware basados en
variaciones físicas del silicio.

El núcleo PLPUF fue implementado exitosamente en Verilog con ancho
parametrizable, configurado a 128~bits, empleando una cadena de inversores con
retroalimentación tipo Fibonacci y una máquina de estados finitos de cuatro
estados. La síntesis en AMD Vivado generó el \textit{bitstream} para la FPGA
Artix-7, y las pruebas de validación funcional confirmaron que el módulo
genera respuestas de 128~bits que varían con el desafío y la duración de
activación, cumpliendo así el primer objetivo específico. El uso de los
atributos de síntesis \texttt{DONT\_TOUCH} y
\texttt{ALLOW\_COMBINATORIAL\_LOOPS} resultó esencial para preservar la
integridad del lazo combinacional durante la optimización del sintetizador.

La interfaz AXI4-Lite \textit{slave}, con un mapa de 12~registros de 32~bits
para control, duración, desafío, respuesta y estado, operó correctamente en
las 10~corridas de evaluación ejecutadas. La verificación de lectura y
escritura de todos los registros, incluida la condición de solo lectura de los
registros de respuesta, validó la integridad del protocolo de comunicación
entre el procesador y el periférico, satisfaciendo el segundo objetivo
específico.

La integración completa del SoC en Vivado IP~Integrator, incluyendo
MicroBlaze~V, memoria BRAM de 64\,KB, controlador UART, módulo de depuración
MDM, generador de reloj y el periférico PLPUF, demostró ser estable y
funcional. El desarrollo del \textit{firmware} de control en C mediante AMD
Vitis, junto con el \textit{driver} de bajo nivel y las dos aplicaciones de
prueba, permitió la ejecución autónoma de más de 140\,000~evaluaciones del
PLPUF distribuidas en 10~corridas independientes, con comunicación estable a
través de la UART. Este resultado confirma el cumplimiento del tercer
objetivo específico.

La evaluación del PLPUF mediante un protocolo exhaustivo basado en las
métricas estándar de la
literatura~\cite{hori_2011_pseudolfsr, maiti_2012_systematic} demostró que el
módulo exhibe propiedades estadísticas coherentes con las esperadas para esta
arquitectura. La uniformidad se mantuvo entre 47\,\% y 51\,\% para todas las
duraciones de activación, indicando ausencia de sesgo estructural
significativo~\cite{maiti_2012_systematic}. La confiabilidad superó el 95\,\%
para duraciones $c \leq 4$, con un máximo de 98.97\,\% para $c = 1$, valor
comparable al reportado por Hori~et~al.~\cite{hori_2011_pseudolfsr} para
FPGAs Virtex-5. La distancia de Hamming inter-desafío alcanzó el 48.48\,\%
para $c = 1$, cercana al ideal de 50\,\%, lo que indica que el PLPUF genera
respuestas suficientemente diversas ante desafíos distintos. Estos resultados
validan la correcta portabilidad de la arquitectura a la familia Artix-7 y
satisfacen el cuarto objetivo específico.

Los datos experimentales confirmaron que la duración de activación~$c$
constituye el parámetro más crítico del PLPUF, al gobernar el compromiso
fundamental entre confiabilidad y entropía. La configuración $c = 1$ ofrece la
mejor combinación de métricas, mientras que valores de~$c \geq 5$ reducen la
confiabilidad por debajo del umbral comúnmente
aceptado~\cite{maiti_2012_systematic}. Esta observación es consistente con los
riesgos documentados por Alsharkawy~et~al.~\cite{alsharkawy_2024_plpufs}
respecto a duraciones elevadas.

El efecto avalancha promedio de 8.43\,\% se situó significativamente por
debajo del ideal de 50\,\%, lo cual constituye una limitación inherente a la
topología de retroalimentación Fibonacci del
PLPUF~\cite{alsharkawy_2024_plpufs}. Esta restricción debe considerarse en
el diseño de protocolos criptográficos que dependan de alta difusión entre el
desafío y la respuesta.

La alta reproducibilidad entre las 10~corridas independientes, con
desviaciones estándar inferiores a un punto porcentual en todas las métricas,
valida tanto la estabilidad del módulo hardware como la robustez del protocolo
de evaluación implementado. No obstante, la evaluación se limitó a un único
dispositivo FPGA, por lo que no fue posible calcular la métrica de unicidad
inter-dispositivo~\cite{wilde_2020_confidence}. La distancia de Hamming
inter-desafío se utilizó como indicador indirecto de la capacidad de
diferenciación del PLPUF.

En síntesis, el proyecto cumplió con todos los objetivos propuestos y generó
una plataforma completa, desde el diseño HDL hasta el software de evaluación
y el \textit{dataset} de pares desafío-respuesta, que puede servir como base
para futuras investigaciones en seguridad de hardware sobre FPGAs en entornos
académicos y de investigación.

% ---------------------------------------------------------------------------
\section{Recomendaciones}

Se recomienda replicar las pruebas sobre múltiples placas Basys\,3, u otras
placas con FPGA Artix-7, para calcular la unicidad inter-dispositivo y el
alias de bit con intervalos de confianza estadísticamente robustos, tal como
sugieren Wilde y Pehl~\cite{wilde_2020_confidence}. Esta evaluación
permitiría cuantificar la capacidad real del PLPUF para distinguir chips
individuales en escenarios de autenticación.

La evaluación realizada se ejecutó a temperatura ambiente y voltaje nominal.
Para validar la confiabilidad del PLPUF en aplicaciones prácticas, es
necesario caracterizar su comportamiento bajo variaciones de temperatura y
fluctuaciones de voltaje, siguiendo las recomendaciones de
Maiti~et~al.~\cite{maiti_2012_systematic}.

Dado que la confiabilidad se degrada progresivamente para duraciones de
activación $c \geq 2$, la incorporación de algoritmos de datos auxiliares
(\textit{Helper Data Algorithms}) con códigos correctores de
errores~\cite{delvaux_2015_helper} permitiría derivar claves criptográficas
determinísticas a partir de las respuestas del PLPUF, incluso en
configuraciones con BER moderado.

El bajo efecto avalancha observado está directamente relacionado con la
estructura del polinomio primitivo utilizado, que posee solo tres puntos de
retroalimentación XOR. Se recomienda evaluar polinomios con mayor número de
taps o funciones de retroalimentación no lineales, como propusieron
Huang~et~al.~\cite{huang_2024_nlfsr}, para mejorar la difusión entre desafío
y respuesta.

Dado que el diseño es parametrizable y está empaquetado como núcleo IP, su
portabilidad a familias FPGA de mayor capacidad, como Kintex-7 o UltraScale+,
permitiría evaluar la influencia de la tecnología de fabricación sobre las
métricas del PLPUF. Asimismo, resulta de interés explorar la compatibilidad
con procesadores suaves alternativos al MicroBlaze~V.

Finalmente, el módulo PLPUF y la infraestructura SoC desarrollados proveen la
base necesaria para implementar protocolos de autenticación ligeros basados en
PUFs, como los propuestos por Idriss~et~al.~\cite{idriss_2021_lightweight} y
Hou~et~al.~\cite{hou_2023_modeling}. Se recomienda utilizar la plataforma
como banco de pruebas para evaluar la viabilidad práctica de estos protocolos
en sistemas embebidos con recursos limitados.

%%% Local Variables:
%%% mode: latex
%%% TeX-master: "main"
%%% End:
